% ============================================================================
% LANGENTWURF LE-02a: La familia - Grundwortschatz
% Spanisch Klasse 7 - Sequenz 2: Mi familia
% ============================================================================
% Tier: 1 (vollständig)
% Doppelstunde: 1-2 von 8
% Fachfarbe: #c41e3a (Karminrot)
% ============================================================================

\documentclass[11pt,a4paper]{article}

\usepackage[utf8]{inputenc}
\usepackage[T1]{fontenc}
\usepackage[ngerman]{babel}
\usepackage{geometry}
\usepackage{fancyhdr}
\usepackage{lastpage}
\usepackage{xcolor}
\usepackage{tcolorbox}
\usepackage{enumitem}
\usepackage{tabularx}
\usepackage{longtable}
\usepackage{array}
\usepackage{booktabs}
\usepackage{amssymb}

% ============================================================================
% KONFIGURATION
% ============================================================================

\newcommand{\fach}{Spanisch}
\newcommand{\fachfarbe}{c41e3a}
\newcommand{\jahrgangsstufe}{7}
\newcommand{\schulart}{Gesamtschule}
\newcommand{\stundenthema}{La familia -- Grundwortschatz}
\newcommand{\reihenthema}{Mi familia}
\newcommand{\stundennummer}{2a}
\newcommand{\reihenstunden}{4}
\newcommand{\gerniveau}{A1}
\newcommand{\lehrwerk}{Que pasa 1}
\newcommand{\lehrwerkseite}{S. 54-57}
\newcommand{\lerngruppengroesse}{24}

% ============================================================================
% LAYOUT
% ============================================================================
\geometry{left=2.5cm, right=2.5cm, top=2.5cm, bottom=2.5cm}
\definecolor{fach}{HTML}{\fachfarbe}
\definecolor{gray}{HTML}{666666}
\definecolor{lightgray}{HTML}{f5f5f5}
\definecolor{successgreen}{HTML}{2a7a4b}
\definecolor{warningorange}{HTML}{b35c00}
\definecolor{basis}{HTML}{2a7a4b}
\definecolor{standard}{HTML}{2c5aa0}
\definecolor{erweiterung}{HTML}{7b1fa2}

\tcbuselibrary{skins,breakable}

\newtcolorbox{infobox}[1][]{
    colback=lightgray,
    colframe=fach,
    fonttitle=\bfseries,
    title=#1,
    breakable
}

\newtcolorbox{scorebox}{
    colback=white,
    colframe=fach,
    fonttitle=\bfseries,
    title=Didaktik-Score,
    breakable
}

\pagestyle{fancy}
\fancyhf{}
\fancyhead[L]{\small\textcolor{gray}{\fach{} -- Jg. \jahrgangsstufe{} -- \stundenthema}}
\fancyhead[R]{\small\textcolor{gray}{LE-\stundennummer{} | Tier 1}}
\fancyfoot[C]{\small\thepage{} / \pageref{LastPage}}
\renewcommand{\headrulewidth}{0.4pt}

% Differenzierungsmarker
\newcommand{\B}{\textcolor{basis}{\textbf{[B]}}}
\newcommand{\Std}{\textcolor{standard}{\textbf{[S]}}}
\newcommand{\E}{\textcolor{erweiterung}{\textbf{[E]}}}

% ============================================================================
% DOKUMENT
% ============================================================================
\begin{document}

% ============================================================================
% DECKBLATT
% ============================================================================
\begin{titlepage}
    \centering
    \vspace*{2cm}
    {\large\textcolor{gray}{\schulart}}\\[2cm]
    {\Huge\textcolor{fach}{\textbf{Langentwurf}}}\\[0.5cm]
    {\large LE-\stundennummer{} | Tier 1 (vollständig)}\\[2cm]
    {\LARGE\textbf{\stundenthema}}\\[0.5cm]
    {\large Sequenz: \reihenthema}\\[2cm]
    \begin{tabular}{rl}
        \textbf{Fach:} & \fach{} (\gerniveau) \\[0.3cm]
        \textbf{Klasse:} & \jahrgangsstufe \\[0.3cm]
        \textbf{Lehrwerk:} & \lehrwerk, \lehrwerkseite \\[0.3cm]
        \textbf{Doppelstunde:} & \stundennummer{} (Std. 1-2 von 8) \\[0.3cm]
    \end{tabular}
    \vfill
\end{titlepage}

\tableofcontents
\newpage

% ============================================================================
% 1. BEDINGUNGSANALYSE
% ============================================================================
\section{Bedingungsanalyse}

\subsection{Lerngruppe}
\begin{tabular}{ll}
    Kursgröße: & \lerngruppengroesse{} SuS \\
    Jahrgangsstufe: & \jahrgangsstufe \\
    Sprachniveau: & \gerniveau{} (GER) -- Anfänger \\
    Fremdsprache: & 2. Fremdsprache (nach Englisch) \\
\end{tabular}

\subsection{Vorwissen aus Sequenz 1}
\begin{itemize}
    \item Begrüßungen und Verabschiedungen
    \item Sich vorstellen: Me llamo..., Soy...
    \item Verb \textit{ser} (yo, tú, él/ella, nosotros)
    \item Herkunft erfragen und angeben: ¿De dónde eres?
    \item Bestimmte und unbestimmte Artikel
\end{itemize}

\subsection{Differenzierung (intern)}
\begin{tabular}{|l|p{3.5cm}|p{3.5cm}|p{3.5cm}|}
\hline
& \B{} \textbf{Basis} & \Std{} \textbf{Standard} & \E{} \textbf{Erweiterung} \\
\hline
\textbf{Ziel} & BR: Rezeption & MR: Produktion mit Hilfe & GY: Freie Produktion \\
\hline
\textbf{Vokabeln} & 6 Kernbegriffe mit Bild & Alle 10 Begriffe & + Großeltern, Cousins \\
\hline
\textbf{Stammbaum} & Beschriften mit Vorgabe & Selbstständig beschriften & Eigene Familie beschreiben \\
\hline
\end{tabular}

\vspace{0.5em}
\textbf{WICHTIG:} Diese Marker sind NUR für die Lehrkraft!

% ============================================================================
% 2. SACHANALYSE
% ============================================================================
\section{Sachanalyse}

\subsection{Wortschatz: Familienmitglieder}

\textbf{Kernwortschatz (10 Begriffe):}
\begin{center}
\begin{tabular}{|l|l|l|}
\hline
\textbf{Spanisch} & \textbf{Deutsch} & \textbf{Artikel} \\
\hline
el padre & der Vater & m \\
la madre & die Mutter & f \\
el hermano & der Bruder & m \\
la hermana & die Schwester & f \\
el abuelo & der Großvater & m \\
la abuela & die Großmutter & f \\
el tío & der Onkel & m \\
la tía & die Tante & f \\
el primo & der Cousin & m \\
la prima & die Cousine & f \\
\hline
\end{tabular}
\end{center}

\textbf{Erweiterungswortschatz:}
\begin{itemize}
    \item los padres (die Eltern), los abuelos (die Großeltern)
    \item el hijo / la hija (der Sohn / die Tochter)
    \item la familia (die Familie)
\end{itemize}

\subsection{Sprachliche Strukturen}

\textbf{Possessivbegleiter (Singular):}
\begin{center}
\begin{tabular}{|l|l|l|}
\hline
\textbf{Pronomen} & \textbf{maskulin} & \textbf{feminin} \\
\hline
yo & mi & mi \\
tú & tu & tu \\
él/ella & su & su \\
\hline
\end{tabular}
\end{center}

\textbf{Redemittel:}
\begin{itemize}
    \item \textbf{Frage:} ¿Quién es? -- Wer ist das?
    \item \textbf{Antwort:} Es mi padre. -- Das ist mein Vater.
    \item \textbf{Erweitert:} Mi hermana se llama Ana.
\end{itemize}

\subsection{Kontrastive Betrachtung}

\begin{tabular}{|l|p{5cm}|p{5cm}|}
\hline
& \textbf{Deutsch} & \textbf{Spanisch} \\
\hline
Genus bei Berufen & Unterschied (Arzt/Ärztin) & Oft Endung -o/-a \\
\hline
Possessiv & Unterscheidet Genus & mi/tu/su einheitlich \\
\hline
Artikel & der/die/das & el/la (einfacher!) \\
\hline
\end{tabular}

\textbf{Positiver Transfer:} Viele Familienbegriffe ähneln Englisch (mother/madre, father/padre).

\subsection{Kulturelle Aspekte}
\begin{itemize}
    \item Familienzusammenhalt in hispanischen Kulturen oft sehr eng
    \item Großeltern leben häufiger im selben Haushalt
    \item Familientreffen (reuniones familiares) haben hohen Stellenwert
    \item Zwei Nachnamen: Vatername + Muttername (z.B. García López)
\end{itemize}

% ============================================================================
% 3. DIDAKTISCHE ANALYSE
% ============================================================================
\section{Didaktische Analyse}

\subsection{Stellung in der Sequenz}
\begin{tabular}{|c|l|l|}
\hline
\textbf{Block} & \textbf{Thema} & \textbf{Status} \\
\hline
\textbf{2a} & \textbf{La familia -- Grundwortschatz} & \textbf{diese Stunde} \\
2b & Mi familia -- Possessivbegleiter & folgt \\
2c & Beschreibungen + tener & folgt \\
2d & Familienpräsentation + Test & folgt \\
\hline
\end{tabular}

\subsection{Kompetenzziele (Rahmenplan MV)}
\begin{itemize}
    \item \textbf{Kommunikativ:} Hörverstehen (Familienmitglieder identifizieren), Sprechen (Familie vorstellen)
    \item \textbf{Sprachlich:} Wortschatz (Familienmitglieder), Grammatik (Artikel, Possessivbegleiter Basics)
    \item \textbf{Interkulturell:} Familienstrukturen in spanischsprachigen Ländern
    \item \textbf{Methodisch:} Wortschatzarbeit mit Bildern, Stammbaum-Arbeit
\end{itemize}

\subsection{Legitimation}
\textbf{Rahmenplan Spanisch MV (Kl. 7-10):}
\begin{itemize}
    \item Kompetenzbereich: Kommunikative Kompetenz -- Sprechen
    \item Standard: ``Die SuS können über sich selbst und ihr persönliches Umfeld (Familie, Freunde) Auskunft geben.''
    \item Themenfeld: ``Persönliches Lebensumfeld'' -- Familie und Freunde
\end{itemize}

% ============================================================================
% 4. STUNDENZIELE
% ============================================================================
\section{Stundenziele}

\subsection{Grobziel}
Die SuS erweitern ihre \textbf{lexikalische Kompetenz im Bereich Familie}, indem sie 10 Familienmitglieder auf Spanisch benennen, in einem Stammbaum zuordnen und einfache Sätze zur Familienbeschreibung bilden.

\subsection{Feinziele (operationalisiert)}
\begin{tabular}{|c|p{7.5cm}|c|c|}
\hline
\textbf{Nr.} & \textbf{Die SuS können...} & \textbf{AFB} & \textbf{Diff.} \\
\hline
FZ1 & die 10 Familienbegriffe (padre, madre, hermano, etc.) korrekt \textbf{aussprechen} und dem deutschen Äquivalent \textbf{zuordnen} & I & alle \\
\hline
FZ2 & einen Stammbaum mit vorgegebenen Begriffen \textbf{beschriften} & I & alle \\
\hline
FZ3 & auf die Frage ``¿Quién es?'' mit ``Es mi...'' \textbf{antworten} & II & alle \\
\hline
FZ4 & einen kurzen Dialog über Familienmitglieder mit Redemitteln \textbf{führen} & II & \Std{}/\E{} \\
\hline
FZ5 & ihre eigene Familie mit 3-4 Sätzen \textbf{beschreiben} & III & \E{} \\
\hline
\end{tabular}

\vspace{1em}
\textbf{AFB-Verteilung:} AFB I: 30\% | AFB II: 50\% | AFB III: 20\%

% ============================================================================
% 5. METHODISCHE ANALYSE
% ============================================================================
\section{Methodische Analyse}

\subsection{Einstieg}
\begin{itemize}
    \item \textbf{Methode:} Bildimpuls -- Foto einer spanischen Familie (z.B. aus Lehrwerk)
    \item \textbf{Begründung:} Aktivierung von Vorwissen, Neugier wecken, Alltagsbezug
    \item \textbf{Alternative:} Stammbaum einer fiktiven Familie an der Tafel
\end{itemize}

\subsection{Erarbeitung}
\begin{itemize}
    \item \textbf{Methode:} Wortschatzeinführung mit Bild-Wort-Zuordnung
    \item \textbf{Begründung:} Multisensorische Verankerung (visuell + auditiv)
    \item \textbf{Scaffolding:} Farbcodierung (männlich = blau, weiblich = rot), Wortkarten
\end{itemize}

\subsection{Übungsphasen}
\begin{itemize}
    \item \textbf{Methode:} Chorisches Sprechen $\rightarrow$ Memory-Spiel $\rightarrow$ Stammbaum-Arbeit
    \item \textbf{Progression:} Rezeptiv $\rightarrow$ Reproduktiv $\rightarrow$ Produktiv
    \item \textbf{Differenzierung:} Stammbaum mit/ohne Vorgaben
\end{itemize}

\subsection{Medien}
\begin{itemize}
    \item Tafel/Whiteboard: Stammbaum-Skizze
    \item Lehrwerk: \lehrwerk, \lehrwerkseite
    \item Arbeitsblatt: AB-02a.pdf
    \item Lernpfad: LP-02a.html (Tablets, optional)
    \item Bildkarten: Familienmitglieder (10 Karten)
\end{itemize}

% ============================================================================
% 6. VERLAUFSPLANUNG
% ============================================================================
\section{Verlaufsplanung}

\textbf{1. Stunde (45 Min.)}

\begin{longtable}{|p{1cm}|p{1.5cm}|p{4cm}|p{3cm}|p{1.5cm}|p{2cm}|}
\hline
\textbf{Zeit} & \textbf{Phase} & \textbf{Lehrerhandeln} & \textbf{Schülerhandeln} & \textbf{SF} & \textbf{Material} \\
\hline
\endhead

0--5 & Einstieg & Begrüßung, Familienfoto zeigen: ``¿Quién es?'' & Vermutungen äußern (deutsch) & Plenum & Bild \\
\hline

5--15 & Input & Vokabeln einführen mit Bildkarten, Aussprache vorsprechen & Nachsprechen, notieren & Plenum & Bildkarten \\
\hline

15--25 & Übung 1 & Chorisches Sprechen: ``El padre, la madre...'' & Mitsprechen, Aussprache üben & Plenum & -- \\
\hline

25--35 & Übung 2 & Memory-Spiel erklären, Gruppen einteilen & Memory spielen (Bild-Wort) & GA & Memory-Karten \\
\hline

35--40 & Sicherung & Schnelle Abfrage: Bild zeigen, Wort nennen & Antworten & Plenum & Bildkarten \\
\hline

40--45 & Überleitung & Merkkasten im AB zeigen, Hausaufgabe & AB beginnen & Plenum & AB-02a \\
\hline
\end{longtable}

\vspace{1em}

\textbf{2. Stunde (45 Min.)}

\begin{longtable}{|p{1cm}|p{1.5cm}|p{4cm}|p{3cm}|p{1.5cm}|p{2cm}|}
\hline
\textbf{Zeit} & \textbf{Phase} & \textbf{Lehrerhandeln} & \textbf{Schülerhandeln} & \textbf{SF} & \textbf{Material} \\
\hline
\endhead

0--5 & Aktivier. & Vokabelquiz: ``Wie heißt Mutter?'' & Antworten: la madre & Plenum & -- \\
\hline

5--15 & Anwendung & Stammbaum an Tafel, gemeinsam beschriften & Begriffe zuordnen, AB bearbeiten & Plenum/EA & Tafel, AB \\
\hline

15--25 & Dialog & Redemittel einführen: ``¿Quién es? -- Es mi...'' & Nachsprechen, PA-Dialoge & PA & Karten \\
\hline

25--35 & Festigung & AB Aufgabe 3+4 bearbeiten lassen, unterstützen & AB bearbeiten & EA & AB-02a \\
\hline

35--40 & Sicherung & 2-3 Dialoge vorführen lassen, Feedback & Präsentieren & Plenum & -- \\
\hline

40--45 & Abschluss & Zusammenfassung, HA: Vokabeln lernen, LP-02a & Notieren & Plenum & -- \\
\hline
\end{longtable}

\textbf{Legende:} EA = Einzelarbeit, PA = Partnerarbeit, GA = Gruppenarbeit

% ============================================================================
% 7. ERWARTETE SCHWIERIGKEITEN
% ============================================================================
\section{Erwartete Schwierigkeiten}

\begin{tabular}{|p{4.5cm}|p{8cm}|}
\hline
\textbf{Schwierigkeit} & \textbf{Reaktion} \\
\hline
Aussprache \textit{hermano} [er'mano] & H ist stumm! Vorsprechen, chorisch wiederholen \\
\hline
Verwechslung \textit{primo/prima} & Visualisierung: Endung -o = männlich, -a = weiblich \\
\hline
Artikel vergessen & Vokabeln immer mit Artikel lernen: \textit{el} padre, \textit{la} madre \\
\hline
Possessiv mi/tu/su & In dieser Stunde nur \textit{mi} einführen, Rest in 2b \\
\hline
Zeitproblem & Puffer: Memory kürzen, als HA aufgeben \\
\hline
\end{tabular}

% ============================================================================
% 8. HAUSAUFGABE
% ============================================================================
\section{Hausaufgabe}

\begin{itemize}
    \item Lehrwerk S. 56, Übung 2 (Familienstammbaum beschriften)
    \item Vokabeln lernen: 10 Familienmitglieder mit Artikel
    \item Optional: Lernpfad LP-02a.html bearbeiten
    \item Optional: Eigenen Familienstammbaum zeichnen (für 2b)
\end{itemize}

% ============================================================================
% 9. LÖSUNGEN
% ============================================================================
\section{Lösungen}

\subsection{AB-02a Lösungen}

\textbf{Aufgabe 1: Zuordnung (4P)}
\begin{itemize}
    \item el padre $\rightarrow$ der Vater
    \item la madre $\rightarrow$ die Mutter
    \item el hermano $\rightarrow$ der Bruder
    \item la hermana $\rightarrow$ die Schwester
\end{itemize}

\textbf{Aufgabe 2: Lückentext (4P)}
\begin{enumerate}[label=\alph*)]
    \item Mi \textbf{padre} se llama Carlos.
    \item Mi \textbf{madre} se llama María.
    \item Tengo un \textbf{hermano} y una \textbf{hermana}.
    \item Mi \textbf{abuelo} es muy simpático.
\end{enumerate}

\textbf{Aufgabe 3: Stammbaum (6P)}
\begin{itemize}
    \item 1 = el abuelo, 2 = la abuela
    \item 3 = el padre, 4 = la madre
    \item 5 = el tío, 6 = la tía
    \item 7 = yo, 8 = el hermano / la hermana
    \item 9 = el primo, 10 = la prima
\end{itemize}

\textbf{Aufgabe 4: Mini-Dialog (6P)}
\begin{itemize}
    \item A: ¿Quién es? (zeigt auf Bild)
    \item B: Es mi padre. Se llama Pedro.
    \item A: ¿Y quién es ella?
    \item B: Es mi madre. Se llama Ana.
\end{itemize}

\textit{Bewertung: Korrekte Struktur (2P), Familienbegriff (2P), Name (1P), Aussprache (1P)}

% ============================================================================
% 10. ANLAGEN
% ============================================================================
\section{Anlagen}

\begin{enumerate}
    \item Arbeitsblatt (AB-02a.pdf)
    \item Musterlösung (ML-02a.pdf)
    \item Lehrerhinweise (LH-02a.pdf)
    \item Lernpfad (LP-02a.html)
    \item Bildkarten Familienmitglieder (10 Stück)
    \item Memory-Karten (Kopiervorlage)
\end{enumerate}

% ============================================================================
% 11. DIDAKTIK-SCORE
% ============================================================================
\newpage
\section{Didaktik-Score}

\begin{scorebox}
\textbf{Bewertung der didaktischen Qualität nach 10 Dimensionen}\\
\textbf{Ziel-Score: $\geq$ 9.0 (Premium-Qualität)}

\vspace{1em}
\begin{tabular}{|c|l|c|c|p{4cm}|}
\hline
\textbf{\#} & \textbf{Dimension} & \textbf{Gew.} & \textbf{Score} & \textbf{Notiz} \\
\hline
1 & Differenzierung & 15\% & 9/10 & [B]/[S]/[E] qualitativ verschieden \\
\hline
2 & Taxonomie (AFB) & 10\% & 9/10 & 30/50/20 gut verteilt \\
\hline
3 & Lernziel-Alignment & 10\% & 10/10 & RP MV explizit adressiert \\
\hline
4 & Aktivierung & 10\% & 10/10 & Memory, Dialoge, Stammbaum \\
\hline
5 & Scaffolding & 10\% & 9/10 & Gestufte Hilfen, Farbcodierung \\
\hline
6 & Feedback-Qualität & 10\% & 9/10 & LP: elaboriertes Feedback \\
\hline
7 & Sprachliche Klarheit & 10\% & 10/10 & Kl. 7 angemessen \\
\hline
8 & Motivation & 10\% & 9/10 & Alltagsbezug Familie, Spiele \\
\hline
9 & Inklusion (UDL) & 10\% & 9/10 & Visuell + auditiv + kinästhetisch \\
\hline
10 & Praktikabilität & 5\% & 9/10 & 90 Min. realistisch geplant \\
\hline
\multicolumn{3}{|r|}{\textbf{GESAMT:}} & \textbf{9.3/10} & \\
\hline
\end{tabular}

\vspace{1em}
$\checkmark$ \textcolor{successgreen}{\textbf{FREIGABE} (Score $\geq$ 9.0)}
\end{scorebox}

\end{document}
